\documentclass[10pt,a4paper]{amsart}
\usepackage[margin=19mm]{geometry}
\usepackage{amsmath,amssymb,mathtools}
\usepackage[scr=rsfs,cal=euler]{mathalfa}
\usepackage{xfrac}
\usepackage[unicode,hidelinks,bookmarks=false]{hyperref}
\newcommand{\ie}{i.e.\ }
\newcommand{\cf}{cf.\ }
\newcommand{\resp}{respectively\ }
\newcommand{\Subsection}{\subsection}
\providecommand{\nobreakdash}{\nobreak\hbox{-}}
\setlength{\emergencystretch}{2em}
\multlinegap0pt
\usepackage{bm}
\usepackage{bbm}
\usepackage{dsfont}
\usepackage{xspace}
\usepackage{cancel}
\usepackage{mathtools}
\usepackage{amsthm}
\makeatletter
\@ifundefined{definition}{\newtheorem{definition}{Definition}}{}
\@ifundefined{lemma}{\newtheorem{lemma}[definition]{Lemma}}{}
\@ifundefined{proposition}{\newtheorem{proposition}[definition]{Proposition}}{}
\makeatother
\newcommand{\B}{\mathbf{B}}
\newcommand{\C}{\mathbb{C}}
\newcommand{\R}{\mathbb{R}}
\newcommand{\dd}{{\rm d}}
\newcommand{\e}{\epsilon}
\newcommand{\proj}{\mathbf{Proj}}
\newcommand{\p}{\partial}
\newcommand{\sph}{{\mathbf{S}^{d-1}}}
\title[Expected signature of stopped Brownian motion on $d$-dimensi]{Expected signature of stopped Brownian motion on $d$-dimensional $C^{2, \alpha}$-domains has finite radius of convergence everywhere: $2\leq d \leq 8$}
\author{Siran Li, Hao Ni}
\date{Source: arXiv:2011.07917v4 (CC BY 4.0)}
\begin{document}
\maketitle

\noindent\emph{Source and licence.} Expected signature of stopped Brownian motion on $d$-dimensional $C^{2, \alpha}$-domains has finite radius of convergence everywhere: $2\leq d \leq 8$. Version arXiv:2011.07917v4 (CC BY 4.0), distributed under the Creative Commons Attribution 4.0 International licence, as recorded in the arXiv licence metadata for this exact version and shown on the abstract page https://arxiv.org/abs/2011.07917v4. Full rights notice: licenses/arxiv-2011.07917-CC-BY-4.0.txt.\par\smallskip

\noindent\textit{Editorial scope.} Counted unit: the Proposition whose source label is \texttt{propn: ODE for hyperbolic dev in d-dim} in the article (source0.tex lines 1406--1417, source label \texttt{propn: ODE for hyperbolic dev in d-dim}), delivered together with the article's own complete original proof of it (source0.tex lines 1419--1485, one \texttt{proof} environment of 67 lines). No mathematics is written, repaired or re-derived: statement and proof are the article's own text line for line. Required context retained uncounted: Lemma\_F\_bar (lines 797--803); lem: sparseness (lines 1349--1351); lem: PDE\_H\_d\_ball (lines 1394--1400); new, B (lines 1511--1514). Every definition, display and label of those lines is retained unchanged. Omitted from this package and not redistributed: 1--796, 804--1348, 1352--1393, 1401--1405, 1418--1418, 1486--1510, 1515--2257. Nothing inside a retained excerpt is omitted or altered: every retained body is delivered exactly as the article's lines are written, so no omission receipt of any kind accompanies this package. Citations: the retained excerpts carry no citation command at all, so no bibliography entry of the article is transcribed and no reference is redistributed. Reference adaptation: none was needed and none was made; every reference inside the delivered unit resolves inside it, so no unresolved reference or citation is produced. Printed slips kept literal: no printed word, symbol, hypothesis or number of the counted statement or of its proof was altered. Layout and class: the article's own class is \texttt{amsart} with packages including amsmath, amssymb, amsfonts, epsfig, mathrsfs, fontenc, color, array, amsthm, amstext, graphicx, setspace, geometry, bbm; the wrapper is the lane's pinned \texttt{amsart} editorial wrapper at 10pt on A4, which restates the theorem-like environments the retained text needs and adds the article's own macro definitions that the retained text needs (\texttt{\textbackslash B}, \texttt{\textbackslash C}, \texttt{\textbackslash R}, \texttt{\textbackslash dd}, \texttt{\textbackslash e}, \texttt{\textbackslash p}, \texttt{\textbackslash proj}, \texttt{\textbackslash sph}), with extra wrapper packages (bm, bbm, dsfont, xspace, cancel, mathtools). The delivered numbering is the unit's own. Assets: no retained span contains a figure environment, a graphics-inclusion command, a PDF-page inclusion, a table or a TikZ picture; no image file is copied into this package, the package declares no asset dependency and no image file is redistributed with it. Licence: version arXiv:2011.07917v4 of this article is distributed under the Creative Commons Attribution 4.0 International licence, as recorded in the arXiv licence metadata for this exact version and shown on the abstract page https://arxiv.org/abs/2011.07917v4. A full rights notice is delivered at licenses/arxiv-2011.07917-CC-BY-4.0.txt. Characters: every retained excerpt is pure ASCII, so the delivered engine, XeLaTeX with its default Latin Modern text font, reports no missing character.
\begin{lemma}[Separation of variables]\label{Lemma_F_bar}
For any $z \in \B(0,\e) \subset \Omega$ and $R \in SO(d)$, it holds that
\begin{equation}\label{eqn_F_bar}
\overline{\mathcal{F}_{\lambda,\e}}(z) =
(R \oplus {\bf id})   \overline{{\mathcal{F}}_{\lambda,\e}}\left( (r, \underbrace{0, \cdots, 0}_{d-1 \text{ zeros}})^{\top}\right).
\end{equation}
\end{lemma}

\begin{lemma}[Sparseness of $\mathcal{H}_{\lambda, \mathbf{D}}$]\label{lem: sparseness}
For each $\lambda >0$, the hyperbolic development $\mathcal{H}_{\lambda, \mathbf{D}}\left(\mathbf{r}\right) = \left (\mathcal{H}^{(1)}_{\lambda, \mathbf{D}}(\mathbf{r}), \cdots, \mathcal{H}^{(d+1)}_{\lambda, \mathbf{D}}(\mathbf{r})\right )^{\top}$ evaluated at the point ${\bf r}=(r,0,\cdots,0)^\top\in \mathbb{R}^{d}$ satisfies $\mathcal{H}^{(i)}_{\lambda, \mathbf{D}}(\mathbf{r})=0$ for $i \in \{2, 3, \cdots, d\}$ and $\mathcal{H}^{(d+1)}_{\lambda, \mathbf{D}}(\mathbf{r}) \geq 1$ almost surely, as long as it is well-defined for such $\lambda$.
\end{lemma}

\begin{lemma}\label{lem: PDE_H_d_ball}
There exists $\lambda^* >0$ such that for any $\lambda \in \C$ with $|\lambda| < \lambda^*$, $\mathcal{H}_{\lambda, \mathbf{D}}$ is $C^2$ on $\mathbf{D}$. Moreover, it satisfies
\begin{equation*}
\Delta \mathcal{H}_{\lambda, \mathbf{D}}(z)=-2\lambda\sum_{i=1}^{d}He_{i}\frac{\partial \mathcal{H}_{\lambda, \mathbf{D}}}{\partial z^{i}}(z)-\lambda^{2}\left(\sum_{i=1}^{d}\left(He_{i}\right)^{2}\right)\mathcal{H}_{\lambda, \mathbf{D}}(z)
\end{equation*}
for each $z \in \mathbf{D}$, and it is subject to the boundary condition $\mathcal{H}_{\lambda, \mathbf{D}} = [0,\ldots, 0, 1]^{\top}$ on $\partial \mathbf{D}$. This PDE boundary value problem has a unique solution.
\end{lemma}

\begin{align}\label{new, B}
    \overline{\mathcal{H}_{\lambda, \mathbf{D}}}(z)  = \left[
\frac{z^{1}}{r}h_{ \lambda}^{(1)}(r),\frac{z^{2}}{r}h_{ \lambda}^{(1)}(r),\, \cdots ,\,\frac{z^{d}}{r}h_{ \lambda}^{(1)}(r),\, h^{(d+1)}_{\lambda}(r) \right]^\top,
\end{align}

\begin{proposition}[ODE system for $\mathcal{H}_{\lambda, \mathbf{D}}$]\label{propn: ODE for hyperbolic dev in d-dim}
Write $h_{\lambda}^{(1)}:=\overline{\mathcal{H}^{(1)}_{\lambda, \mathbf{D}}}$ and $h_{\lambda}^{(d+1)}:=\overline{\mathcal{H}^{(d+1)}_{\lambda, \mathbf{D}}}$ where 
\begin{equation*}
    \overline{\mathcal{H}_{\lambda, \mathbf{D}}}(r) := \int_{SO(d)} \mathcal{H}_{\lambda, \mathbf{D}}\left(R\cdot\left[r,\underbrace{0,\ldots,0}_{d-1\,\text{zeros}}\right]^\top\right) \,\dd\chi(R).
\end{equation*}
Then for $d\in \{2,3,4,\ldots\}$ it holds that
  \begin{eqnarray}
    &&\left(h^{(1)}_{\lambda} \right)'' + \frac{d-1}{r} \left(h^{(1)}_{\lambda}\right)'- \frac{(d-1) h^{(1)}_{\lambda}}{r^2} = -2\lambda \left(h^{(d+1)}_{\lambda}\right)' - \lambda^{2}h^{(1)}_{\lambda},\label{ODE one, d-dim}\\
    && \left(h^{(d+1)}_{\lambda}\right)'' +\frac{d-1}{r}\left(h^{(d+1)}_{\lambda}\right)' = - 2 \lambda \left\{\frac{d-1}{r}h^{(1)}_{\lambda} +\left(h^{(1)}_{\lambda}\right)'\right\} -  d \lambda^{2}h^{(d+1)}_{\lambda},\label{ODE two, d-dim}\\
    && h^{(1)}_\lambda (1) = 0,\qquad h^{(d+1)}_\lambda(1)=1.
    \end{eqnarray}
\end{proposition}

\begin{proof} By Lemma~\ref{Lemma_F_bar} we have $
{\mathcal{H}_{\lambda, \mathbf{D}}}(z) = \left(R \oplus \mathbf{id} \right) \overline{\mathcal{H}_{\lambda, \mathbf{D}}}(r)$ for $z =R\cdot(r,0,\ldots,0)^\top$. Also,  Lemma~\ref{lem: sparseness} gives us $
\overline{\mathcal{H}_{\lambda, \mathbf{D}}}(z) \in \left\{ w \in \mathbb{R}^{d+1}:\, w^{(k)} = 0\text{ for } k \in \{2, \cdots, d\} \text{ and }  w^{(d+1)}\geq 1\right\}.$ Thus
\begin{align*}
    {\mathcal{H}_{\lambda, \mathbf{D}}}(z) =\begin{bmatrix}
    R \cdot \begin{bmatrix}
    h_\lambda^{(1)}(r)\\
    0\\
    \vdots\\
    0
    \end{bmatrix} \\
 h_\lambda^{(d+1)}(r)
    \end{bmatrix} = \begin{bmatrix}
\frac{z^{1}}{r}h_{ \lambda}^{(1)}(r) \\\frac{z^{2}}{r}h_{ \lambda}^{(1)}(r)\\ \vdots \\\frac{z^{d}}{r}h_{ \lambda}^{(1)}(r)\\h^{(d+1)}_{\lambda}(r) \end{bmatrix},
\end{align*}
since the unit sphere $\sph$ is invariant under $R$.


Consider now the PDE system in Lemma~\ref{lem: PDE_H_d_ball}:
\begin{equation}\label{new, PDE, C}
\Delta \mathcal{H}_{\lambda, \mathbf{D}}=-2\lambda\sum_{i=1}^{d}He_{i}\frac{\partial \mathcal{H}_{\lambda, \mathbf{D}}}{\partial z^{i}}-\lambda^{2}\left(\sum_{i=1}^{d}\left(He_{i}\right)^{2}\right)\mathcal{H}_{\lambda, \mathbf{D}}.
\end{equation}
The coefficients can be  easily computed: since $H(e_i)$ is the matrix whose only nonzero entries are the $(i, d+1)^{\text{th}}$ and $(d+1, i)^{\text{th}}$ ones (both equal to $1$), we find that
\begin{equation}\label{proj identity A}
    \left(\sum_{i=1}^{d}\left(He_{i}\right)^{2}\right) = {\rm diag}(1,\cdots,1,d).
\end{equation}
Moreover, the $\alpha^{\text{th}}$ component of  $\sum_{i=1}^{d}He_{i}\frac{\partial \mathcal{H}_{\lambda, \mathbf{D}}}{\partial z^{i}}$ equals $
      {\p \mathcal{H}_{\lambda, \mathbf{D}}^{(d+1)}}\slash{\p z^\alpha}$ if $\alpha\in \{1,2,\ldots,d\}$ and $
      \sum_{i=1}^d \p \mathcal{H}_{\lambda, \mathbf{D}}^{(i)}\slash \p z^i$ if  $\alpha=d+1$. In the latter case we further have
      \begin{align}\label{proj identity B}
    \sum_{i=1}^d \frac{\p \mathcal{H}_{\lambda, \mathbf{D}}^{(i)}}{\p z^i}(z) &=   \sum_{i=1}^d\left\{ \frac{1}{r}h_{\lambda}^{(1)}(r) - \frac{z^i z^i}{r^3}h_{\lambda}^{(1)}(r) + \frac{z^i}{r}\left(h_{\lambda}^{(1)}(r)\right)'(r) \frac{\p r}{\p z_i} \right\}\nonumber \\
   &= \frac{d-1}{r}h_{\lambda}^{(1)}(r)+ \left(h_{\lambda}^{(1)}\right)'(r). 
      \end{align}

On the other hand, $z^1/r=\cos\theta_1$ in the usual spherical co-ordinates $\omega=z/r=(\theta_1, \ldots, \theta_{d-1})$ on $\mathbf{S}^{d-1}$. Note also that $\Delta$ on the left-hand side of Eq.~\eqref{new, PDE, C} is the Euclidean Laplacian on $\R^{d+1}=\R^d\times \R$; when acting on the $\R^d$-component, in the spherical polar co-ordinates 
\begin{align}\label{identity, beltrami}
    \Delta = \p_{rr} + \frac{d-1}{r} \p_r  + \frac{1}{r^2}\Delta_{\mathbf{S}^{d-1}},  
\end{align}
where $\Delta_{\mathbf{S}^{d-1}}$ is the Laplace--Beltrami operator on the sphere $\mathbf{S}^{d-1}$.


To proceed, let us project Eq.~\eqref{new, PDE, C} onto the first and the last components. We shall prove the blowup of the resulting ODE system for $h_\lambda^{(1)}$ and $h_\lambda^{(d+1)}$.



For the projection onto the $(d+1)^{\text{th}}$ component, note that $\left[\proj_{d+1} \mathcal{H}_{\lambda,\mathbf{D}}\right](z)=h_\lambda^{(d+1)}(r)$  by Eq.~\eqref{new, B}, which is independent of $\omega \in \mathbf{S}^{d-1}$. So Identity~\eqref{identity, beltrami} leads to
\begin{align*}
    \Delta \proj_{d+1} \mathcal{H}_{\lambda,\mathbf{D}}(z) = \left(h_\lambda^{(d+1)}\right)''(r) + \frac{d-1}{r}\left(h_\lambda^{(d+1)}\right)'(r).
\end{align*}
The projection of the right-hand side of Eq.~\eqref{new, PDE, C} can be found via Eqs.~\eqref{proj identity A} and \eqref{proj identity B}. Thus we obtain Eq.~\eqref{ODE two, d-dim}. 


It remains to prove Eq.~\eqref{ODE one, d-dim}. To this end, note that Eqs.~\eqref{new, PDE, C} -- \eqref{identity, beltrami} yield that
\begin{align*}
    \Delta\left(\frac{z^1}{r} h_\lambda^{(1)}\right) = -2\lambda \frac{\p h_\lambda^{(d+1)}}{\p z^1} - \lambda^2\frac{z^1}{r}  h_1.
\end{align*}
Writing in the spherical polar co-ordinates, we have $z^1\slash r = \cos\theta_1$, $\p h_\lambda^{(d+1)}\slash{\p z^1} = \cos\theta_1 \left( h_\lambda^{(d+1)} \right)'$, as well as 
\begin{align*}
     \Delta\left(\frac{z^1}{r} h_\lambda^{(1)}\right) = \cos\theta_1\left\{ \left(h_\lambda^{(1)}\right)''+\frac{d-1}{r}\left(h_\lambda^{(1)}\right)'\right\}  + \frac{h_\lambda^{(1)}}{r^2} \Delta_{\mathbb{S}^{d-1}} \cos\theta_1.
\end{align*}
In addition, the Laplace--Beltrami $\Delta_{\mathbf{S}^{d-1}}$ on $\mathbf{S}^{d-1}$ can be expressed in terms of the Laplace--Beltrami $\Delta_{\mathbf{S}^{d-2}}$ on $\mathbf{S}^{d-2}$ --- for $\omega = (\theta_1, \omega')$ where $\omega'=(\theta_2,\cdots,\theta_{d-1}) \in \mathbf{S}^{d-2}$, one has
\begin{align*}
    \Delta_{\mathbf{S}^{d-1}} = \left(\sin\theta_1\right)^{2-d} \frac{\p}{\p\theta_1} \left( \left(\sin\theta_1\right)^{d-2}\frac{\p}{\p \theta_1} \right) + \frac{1}{\sin^2 \theta_1} \Delta_{\mathbf{S}^{d-2}},
\end{align*}
from which it follows that $$\Delta_{\mathbf{S}^{d-1}} \cos\theta_1 = -(d-1)\cos\theta_1.$$ We thus obtain Eq.~\eqref{ODE one, d-dim}.

The boundary conditions $h^{(1)}_\lambda (1) = 0$, $h^{(d+1)}_\lambda(1)=1$ follow from $\Phi_{\mathbf{D}}(1)=1$, which holds since the stopping time is zero on the boundary $\p\mathbf{D}$.   \end{proof}

\end{document}
